\documentclass[11pt]{article}
\newcommand{\DocShort}{Guía de estudio - Parcial 1}
\usepackage{fontspec}
\setmainfont{Latin Modern Roman}
\setsansfont{Latin Modern Sans}
\setmonofont{Latin Modern Mono}
\usepackage[spanish,es-noshorthands]{babel}
\decimalpoint
\usepackage{geometry}
\geometry{letterpaper, top=2.5cm, bottom=2.5cm, left=2.8cm, right=2.8cm}
\usepackage[expansion=false]{microtype}
\usepackage{parskip}
\setlength{\parskip}{6pt}
\usepackage{amsmath,amssymb}
\usepackage{booktabs}
\usepackage{array}
\usepackage{longtable}
\usepackage{caption}
\usepackage[dvipsnames,table]{xcolor}
\usepackage{enumitem}
\setlist[itemize]{topsep=2pt, itemsep=2pt, parsep=0pt}
\setlist[enumerate]{topsep=2pt, itemsep=2pt, parsep=0pt}
\usepackage{fancyhdr}
\usepackage[hidelinks,pdfencoding=auto]{hyperref}
\usepackage{titlesec}
\usepackage{tcolorbox}
\tcbuselibrary{breakable}
\usepackage{pgfplots}
\pgfplotsset{compat=1.18}
\usepackage{tikz}
\usepackage{multicol}
\usepackage{makecell}
\usepackage{siunitx}

\sisetup{
  group-separator={,},
  group-minimum-digits=4,
  output-decimal-marker={.}
}

\definecolor{darkblue}{HTML}{1B3A5C}
\definecolor{medblue}{HTML}{2E6DA4}
\definecolor{lightblue}{HTML}{D6E8F7}
\definecolor{teal}{HTML}{1A7A6E}
\definecolor{lightteal}{HTML}{D0EDEA}
\definecolor{olive}{HTML}{4A6A1E}
\definecolor{lightolive}{HTML}{E2EDD0}
\definecolor{lightgrey}{HTML}{F5F5F5}
\definecolor{midgrey}{HTML}{6F6F6F}
\definecolor{warnyellow}{HTML}{FFF3CD}
\definecolor{warnorange}{HTML}{7B3F00}

\titleformat{\section}
{\large\bfseries\color{darkblue}}
{\thesection}{0.75em}{}[{\color{medblue}\titlerule[0.8pt]}]
\titleformat{\subsection}
{\normalsize\bfseries\color{darkblue}}
{\thesubsection}{0.75em}{}
\titleformat{\subsubsection}
{\normalsize\bfseries\color{medblue}}
{\thesubsubsection}{0.75em}{}

\pagestyle{fancy}
\fancyhf{}
\fancyhead[L]{\small\color{midgrey}ICV 342 Hidrología — PUCMM}
\fancyhead[R]{\small\color{midgrey}\DocShort}
\fancyfoot[C]{\small\color{midgrey}\thepage}
\setlength{\headheight}{14pt}
\renewcommand{\headrulewidth}{0.4pt}
\renewcommand{\footrulewidth}{0pt}

\rowcolors{2}{lightgrey}{white}
\renewcommand{\arraystretch}{1.12}

\newtcolorbox{infobox}{
  breakable,
  colback=lightblue,
  colframe=medblue,
  boxrule=0.8pt,
  arc=3pt,
  left=8pt,
  right=8pt,
  top=7pt,
  bottom=7pt
}

\newtcolorbox{warnbox}{
  breakable,
  colback=warnyellow,
  colframe=orange!70!black,
  boxrule=0.8pt,
  arc=3pt,
  left=8pt,
  right=8pt,
  top=7pt,
  bottom=7pt
}

\newtcolorbox{conceptbox}{
  breakable,
  colback=lightteal,
  colframe=teal,
  boxrule=0.8pt,
  arc=3pt,
  left=8pt,
  right=8pt,
  top=7pt,
  bottom=7pt
}

\newtcolorbox{resultbox}{
  breakable,
  colback=lightolive,
  colframe=olive,
  boxrule=0.8pt,
  arc=3pt,
  left=8pt,
  right=8pt,
  top=7pt,
  bottom=7pt
}

\newcommand{\doccover}[2]{
  \begin{center}
  {\small\color{midgrey}Pontificia Universidad Católica Madre y Maestra\\Escuela de Ingeniería Civil y Ambiental}

  \vspace{1.0cm}
  {\color{medblue}\rule{\linewidth}{2pt}}
  \vspace{0.3cm}

  {\Large\bfseries\color{darkblue}ICV 342 Hidrología}\\[0.35em]
  {\LARGE\bfseries\color{darkblue}#1}\\[0.2em]
  {\small\color{midgrey}#2}

  \vspace{0.3cm}
  {\color{medblue}\rule{\linewidth}{2pt}}
  \vspace{0.9cm}
  \end{center}

  \begin{center}
  \begin{tabular}{ll}
  \toprule
  \textbf{Asignatura} & ICV 342 Hidrología \\
  \textbf{Profesor} & Ricardo Rafael Hernández Moreira, PhD \\
  \textbf{Alcance} & Unidades I, II y III \\
  \bottomrule
  \end{tabular}
  \end{center}

  \vspace{0.4cm}
}

\newcommand{\handoutintro}[2]{
  \begin{infobox}
  \textbf{Propósito.} #1

  \vspace{0.55\baselineskip}

  \textbf{Cuándo conviene.} #2
  \end{infobox}
}

\newcommand{\resultnote}[1]{
  \begin{resultbox}
  \textbf{Idea clave.} #1
  \end{resultbox}
}

\newcommand{\conceptnote}[1]{
  \begin{conceptbox}
  \textbf{Concepto clave.} #1
  \end{conceptbox}
}

\newcommand{\warningnote}[1]{
  \begin{warnbox}
  \textbf{Atención.} #1
  \end{warnbox}
}

\AtBeginDocument{
  \hypersetup{
    pdfauthor={Ricardo Rafael Hernández Moreira},
    pdfsubject={ICV 342 Hidrología},
    pdftitle={\DocShort}
  }
}


\begin{document}
\doccover{Guía de estudio para el Parcial 1}{Repaso del curso}

\handoutintro
{Preparar al estudiante para un parcial principalmente de selección múltiple sobre las Unidades I, II y III de Hidrología. La guía está pensada para repaso activo: leer, comparar conceptos, resolver casos breves y detectar errores frecuentes.}
{Conviene usarla como hoja de estudio principal antes del examen y como base para repasar en clase o en tutoría.}

\section{Cómo estudiar este parcial}
\begin{itemize}
\item No memorices solo definiciones aisladas; busca relaciones: ciclo hidrológico, clima, cuenca, disponibilidad y gestión.
\item Para cada tema, responde tres preguntas: qué es, cómo se mide o identifica y por qué importa.
\item Si puedes explicar un concepto con un ejemplo dominicano, lo dominas mejor.
\item La lectura del Plan Hidrológico Nacional debe ser superficial pero correcta: entiende qué aborda cada capítulo relevante, no todos los detalles numéricos.
\end{itemize}

\section{Mapa del parcial}
\begin{center}
\begin{tabular}{p{0.16\linewidth} p{0.34\linewidth} p{0.42\linewidth}}
\toprule
\textbf{Unidad} & \textbf{Qué domina} & \textbf{Qué puede preguntarse} \\
\midrule
I & Introducción, ciclo hidrológico y contexto nacional & Definiciones, problemas del agua en RD, PHN, usos y disponibilidad. \\
II & Meteorología e hidrología & Variables climáticas, estaciones, evaporación, tiempo vs clima, impacto del clima. \\
III & Cuencas hidrográficas & Delimitación, fisiografía, drenaje, respuesta hidrológica, efectos del uso del suelo. \\
\bottomrule
\end{tabular}
\end{center}

\section{Unidad I. Introducción, ciclo hidrológico y contexto nacional}
\begin{conceptbox}
\textbf{Debes poder decir:} la hidrología estudia el agua en la Tierra, su circulación, distribución y propiedades, con interés especial en cómo esa circulación afecta disponibilidad, uso y gestión.
\end{conceptbox}

\subsection{Ideas que no deben faltar}
\begin{itemize}
\item Definición de hidrología.
\item Ciclo hidrológico como sistema de transferencias entre atmósfera, superficie, subsuelo y océanos.
\item Disponibilidad de agua no equivale a abundancia promedio anual.
\item Problemas hidrológicos: desigualdad temporal, desigualdad espacial, calidad, demanda, gestión, eventos extremos.
\item En República Dominicana importa la lectura territorial del recurso, no solo la cantidad total.
\end{itemize}

\subsection{Plan Hidrológico Nacional}
\begin{resultbox}
Memoriza el mapa funcional del PHN:
\begin{itemize}
\item Capítulo 2: marco físico y climatológico.
\item Capítulo 5: caracterización del ámbito del estudio.
\item Capítulo 6: recursos hídricos, demandas y balances.
\end{itemize}
\end{resultbox}

\subsection{Lo que debes explicar}
\begin{itemize}
\item Por qué un país puede tener agua y aun así sufrir escasez.
\item Por qué la gestión del agua incluye cantidad, calidad, tiempo y lugar.
\item Cómo se conecta el PHN con la idea de seguridad hídrica.
\end{itemize}

\subsection{Preguntas de repaso}
\begin{enumerate}
\item Define hidrología en una sola oración.
\item Explica la diferencia entre disponibilidad y demanda.
\item Resume el papel de los capítulos 2, 5 y 6 del PHN.
\end{enumerate}

\section{Unidad II. Meteorología e hidrología}
\begin{conceptbox}
\textbf{Debes poder distinguir:} tiempo atmosférico y clima, y relacionar variables meteorológicas con procesos hidrológicos.
\end{conceptbox}

\subsection{Variables clave}
\begin{center}
\begin{tabular}{p{0.23\linewidth} p{0.33\linewidth} p{0.34\linewidth}}
\toprule
\textbf{Variable} & \textbf{Qué representa} & \textbf{Efecto hidrológico} \\
\midrule
Precipitación & Entrada de agua desde la atmósfera & Alimenta escurrimiento, infiltración y almacenamiento. \\
Temperatura & Nivel térmico del aire & Influye en evaporación y capacidad del aire para retener vapor. \\
Humedad relativa & Cercanía al estado de saturación & Afecta evaporación y transpiración. \\
Viento & Movimiento del aire & Favorece intercambio de vapor con la superficie. \\
Radiación solar & Energía disponible & Controla calentamiento y pérdidas por evaporación. \\
\bottomrule
\end{tabular}
\end{center}

\subsection{Puntos de estudio}
\begin{itemize}
\item Diferencia entre tiempo y clima.
\item Importancia de la calidad y continuidad de los datos meteorológicos.
\item Evaporación como pérdida de agua hacia la atmósfera.
\item Cambio climático como alteración de patrones de precipitación, temperatura y extremos.
\item Relación entre observación meteorológica y análisis hidrológico.
\end{itemize}

\subsection{Lo que debes explicar}
\begin{itemize}
\item Qué instrumento mide cada variable básica.
\item Cómo una subida de temperatura puede aumentar la evaporación.
\item Por qué los datos meteorológicos son esenciales para planificar recursos hídricos.
\end{itemize}

\subsection{Preguntas de repaso}
\begin{enumerate}
\item ¿En qué se diferencia clima de tiempo atmosférico?
\item Nombra tres variables meteorológicas relevantes para hidrología.
\item Explica por qué la evaporación es una pérdida hidrológica.
\end{enumerate}

\section{Unidad III. Cuencas hidrográficas}
\begin{conceptbox}
\textbf{Debes poder explicar:} la cuenca es la unidad natural de análisis porque organiza el ingreso, movimiento y salida del agua en un territorio.
\end{conceptbox}

\subsection{Ideas centrales}
\begin{itemize}
\item Definición de cuenca hidrográfica.
\item Divisores de agua y delimitación topográfica.
\item Área, forma, relieve y red de drenaje como controladores de la respuesta hidrológica.
\item Densidad de drenaje y orden de corrientes.
\item Parámetros morfométricos como compacidad, elongación y factor de forma.
\item Diferencia entre cuenca superficial y cuenca subterránea.
\item Papel de la cobertura forestal y del uso del suelo.
\end{itemize}

\subsection{Parámetros morfométricos que debes reconocer}
\begin{center}
\begin{tabular}{p{0.26\linewidth} p{0.30\linewidth} p{0.32\linewidth}}
\toprule
\textbf{Parámetro} & \textbf{Qué resume} & \textbf{Qué sugiere} \\
\midrule
Factor de forma & Relación entre área y longitud principal & Qué tan concentrada puede ser la respuesta. \\
Coeficiente de compacidad & Qué tan cerca está la cuenca de una forma compacta & Comparación geométrica entre cuencas. \\
Índice de elongación & Qué tan alargada es la cuenca & Tendencia a repartir o concentrar el escurrimiento. \\
Densidad de drenaje & Longitud total de cauces por área & Rapidez potencial de evacuación. \\
\bottomrule
\end{tabular}
\end{center}

\subsection{Relaciones que debes dominar}
\begin{resultbox}
Si la cuenca es más compacta, con pendientes mayores y red de drenaje eficiente, la respuesta suele ser más rápida y el pico del hidrograma puede ser mayor.
\end{resultbox}

\begin{itemize}
\item La forma de la cuenca afecta el tiempo al pico.
\item La urbanización tiende a aumentar la rapidez de respuesta.
\item Más drenaje concentrado suele implicar menor tiempo de concentración.
\item El bosque, bien conservado, favorece infiltración y control de escorrentía.
\item El factor de forma y el coeficiente de compacidad ayudan a anticipar la concentración del flujo.
\end{itemize}

\subsection{Fórmulas básicas a reconocer}
\begin{align*}
Q &= A \cdot v \\
\bar{P} &= \text{promedio areal de precipitación}
\end{align*}

\begin{align*}
D_d &= \frac{L_t}{A} \\
F_f &= \frac{A}{L^2} \\
C_c &= \frac{P}{2\sqrt{\pi A}}
\end{align*}

\subsection{Lo que debes explicar}
\begin{itemize}
\item Cómo se delimita una cuenca.
\item Qué significa drenaje denso.
\item Cómo el uso del suelo cambia la respuesta de la cuenca.
\end{itemize}

\subsection{Preguntas de repaso}
\begin{enumerate}
\item Define cuenca hidrográfica.
\item Explica qué es el divisor topográfico.
\item Relaciona forma de la cuenca y gasto pico.
\end{enumerate}

\section{Preguntas para practicar}
\begin{enumerate}
\item Diferencia entre tiempo y clima.
\item Relación entre temperatura y evaporación.
\item Lectura del PHN por capítulos 2, 5 y 6.
\item Diferencia entre cuenca superficial y subterránea.
\item Efecto de la urbanización sobre el hidrograma.
\item Interpretación de la densidad de drenaje.
\item Cálculo básico de factor de forma o densidad de drenaje.
\item Función del bosque en la cuenca.
\end{enumerate}

\section{Errores comunes}
\begin{warnbox}
\begin{itemize}
\item Confundir disponibilidad promedio con disponibilidad efectiva en sitio y momento.
\item Tratar clima y tiempo atmosférico como sinónimos exactos.
\item Pensar que el PHN es solo una lista de obras y no un marco de planificación.
\item Creer que cualquier cuenca responde igual ante una tormenta.
\item Olvidar que la calidad de datos meteorológicos afecta el análisis hidrológico.
\end{itemize}
\end{warnbox}

\section{Lista final de repaso}
\begin{itemize}
\item Puedo definir hidrología sin leer.
\item Puedo explicar el ciclo hidrológico con mis palabras.
\item Puedo distinguir tiempo, clima y meteorología.
\item Puedo nombrar capítulos 2, 5 y 6 del PHN y su función.
\item Puedo definir cuenca y explicar por qué importa.
\item Puedo interpretar factor de forma, compacidad y densidad de drenaje.
\item Puedo decir cómo la forma y el drenaje de la cuenca afectan el escurrimiento.
\item Puedo relacionar variables meteorológicas con evaporación y lluvia.
\end{itemize}

\section{Fuentes base}
\begin{itemize}
\item Aparicio-Mijares, \emph{Fundamentos de hidrología de superficie}.
\item Chow, Ven Te, \emph{Hidrología aplicada}.
\item Linsley, Kohler y Paulhus, \emph{Hidrología para ingenieros}.
\item Plan Hidrológico Nacional 2025-2045 de la República Dominicana.
\end{itemize}

\end{document}
